\pdfinfoomitdate=1
\pdftrailerid{}
\pdfsuppressptexinfo=15
\documentclass[11pt]{article}
\usepackage[T1]{fontenc}
\usepackage{lmodern}
\usepackage{amsmath,amssymb,amsthm}
\usepackage{booktabs}
\usepackage{geometry}
\usepackage[hidelinks]{hyperref}
\usepackage{microtype}
\geometry{margin=1in}

\newtheorem{theorem}{Theorem}[section]
\newtheorem{lemma}[theorem]{Lemma}
\newtheorem{proposition}[theorem]{Proposition}
\newtheorem{corollary}[theorem]{Corollary}
\theoremstyle{definition}
\newtheorem{definition}[theorem]{Definition}
\newtheorem{computation}[theorem]{Verified computation}
\theoremstyle{remark}
\newtheorem{remark}[theorem]{Remark}

\newcommand{\QQ}{\mathbf Q}
\newcommand{\PP}{\mathbf P}
\newcommand{\Spec}{\operatorname{Spec}}
\newcommand{\Proj}{\operatorname{Proj}}
\newcommand{\Der}{\operatorname{Der}}
\newcommand{\MC}{\operatorname{MC}}
\newcommand{\im}{\operatorname{im}}
\newcommand{\rank}{\operatorname{rank}}
\newcommand{\ad}{\operatorname{ad}}
\newcommand{\gr}{\operatorname{gr}}
\newcommand{\Tor}{\operatorname{Tor}}

\title{A strict dg Lie proof of the Gr\"unbaum--Sreedharan\
Stanley--Reisner smoothing theorem}
\author{}
\date{}

\begin{document}
\maketitle

\begin{abstract}
We prove that the Stanley--Reisner threefold associated with the
Gr\"unbaum--Sreedharan eight-vertex triangulated three-sphere is the special
fibre of a flat projective scheme over $\QQ[[s]]$ whose geometric generic
fibre is smooth.  The deformation controller is the strict dg Lie algebra of
internal-weight-zero derivations of a Tate resolution relative to the fixed
polynomial ring.  An exact deformation over $\QQ[s]/(s^4)$ is lifted directly
to a square-zero Tate differential.  The ordinary Maurer--Cartan equation,
the strict Bianchi identity, and a certified exact sequence in degrees
$1,2,3$ then extend its prescribed two-jet to all orders.  No transferred
structure and no higher multilinear operation is used.
\end{abstract}

\tableofcontents

\section{Statement and exact input}

Put $k=\QQ$, $S=k[x_1,\ldots,x_8]$, and let $\Delta$ have facets
\[
\begin{gathered}
1234,1237,1267,1347,1567,2345,2367,3467,3456,4567,\\
2358,2368,3568,1268,1568,1248,2458,1478,1578,4578.
\end{gathered}
\]
Its Stanley--Reisner ideal is
\[
\begin{split}
I=(\,&x_6x_7x_8,x_4x_6x_8,x_3x_7x_8,x_3x_5x_7,
x_3x_4x_8,x_2x_7x_8,\\
&x_2x_5x_7,x_2x_5x_6,x_2x_4x_7,x_2x_4x_6,
x_1x_4x_6,x_1x_4x_5,\\
&x_1x_3x_8,x_1x_3x_6,x_1x_3x_5,x_1x_2x_5\,)\subset S.
\end{split}                                                    \tag{1.1}
\]
Write these ordered monomials as $m_1,\ldots,m_{16}$, set $A=S/I$, and put
$X_0=\Proj A\subset\PP^7_k$.  For $n\geq1$ let
\[
 R_n=k[s]/(s^n),\qquad R=k[[s]].
\]

\begin{theorem}[smoothing theorem]\label{thm:main}
There is a flat projective morphism
\[
 X\longrightarrow\Spec\QQ[[s]]
\]
whose special fibre is $X_0$ and whose geometric generic fibre is smooth.
Consequently $X_0$ is smoothable in characteristic zero.
\end{theorem}

This is the exact conclusion proved here.  In particular, no assertion about
the canonical bundle, a Calabi--Yau structure, irreducibility, or a named
smooth threefold is part of the theorem.

\begin{computation}[special fibre]\label{comp:special}
The exact rational computation
\path{../referee-packet/code/verify_combinatorics.py} verifies that the
displayed complex has $20$ tetrahedral facets and precisely the sixteen
minimal nonfaces in (1.1).  It verifies Reisner's vanishing condition on all
$97$ faces by exact boundary-matrix row reduction over $\QQ$.  Thus $A$ is
Cohen--Macaulay of Krull dimension $4$, and $X_0$ is an equidimensional
Cohen--Macaulay projective threefold.  The deterministic output is
\path{../referee-packet/verification/combinatorics_QQ.txt}.
\end{computation}

\section{The strict controller and its classical meaning}

\subsection{Gradings and derivations}

Choose a full internally graded Tate resolution $P\to A$ over $S$ which
continues the independently certified low Tate stages described in
Section~\ref{sec:h3}.  It is a
semi-free graded-commutative dg $S$-algebra in nonnegative homological
degrees, its differential $d$ lowers homological degree by one and preserves
internal weight, and
\[
 H_0(P)=A,\qquad H_i(P)=0\quad(i>0).
\]
Tate resolutions exist by Avramov~\cite[Proposition 6.1.4]{Avramov}; over a
$\QQ$-algebra the usual graded-commutative description agrees with the
divided-power description needed in general.

A cohomological degree-$p$ derivation lowers Tate homological degree by $p$.
Define two strict dg Lie algebras, always retaining internal weight zero,
\[
 \mathfrak g=\Der_k(P,P)_0,
 \qquad
 \mathfrak h=\Der_S(P,P)_0.
\]
For $\theta\in\mathfrak g^p$ and $\eta\in\mathfrak g^q$, set
\[
 \partial\theta=[d,\theta]
   =d\theta-(-1)^p\theta d,
 \qquad
 [\theta,\eta]=\theta\eta-(-1)^{pq}\eta\theta.              \tag{2.1}
\]
The commutator of derivations is again a derivation, and (2.1) gives strict
dg Lie algebras.

The relative algebra $\mathfrak h$ is the correct fixed-coordinate
controller: its degree-zero automorphisms fix $S$, hence fix the eight
projective coordinates.  The absolute algebra $\mathfrak g$ is nevertheless
the one whose $H^1$ is the computed abstract tangent space.  These statements
are compatible, rather than competing, because
\[
 \mathfrak h^p=\mathfrak g^p\quad(p\geq1).                   \tag{2.2}
\]
Indeed, a positive-degree derivation sends $P_0=S$ to a negative homological
degree, which is zero.  It follows that
\[
 H^1(\mathfrak h)\twoheadrightarrow H^1(\mathfrak g),
 \qquad H^q(\mathfrak h)\xrightarrow{\sim}H^q(\mathfrak g)
 \quad(q\geq2).                                              \tag{2.3}
\]
The first arrow need not be an isomorphism.  In fact, the exact package
calculation gives dimensions $109$ and $53$ for the relative embedded
tangent space and the absolute tangent space respectively.  We never identify
them.

The augmentation induces a quasi-isomorphism
\[
 \Der_k(P,P)_0\longrightarrow\Der_k(P,A)_0.                 \tag{2.4}
\]
To see this directly, write the source and target as
$\operatorname{Hom}_P(\Omega_{P/k},P)_0$ and
$\operatorname{Hom}_P(\Omega_{P/k},A)_0$.  The module of differentials of a
semi-free algebra is semi-free, so Hom from it preserves the quasi-isomorphism
$P\to A$; compare Avramov~\cite[Corollary 6.2.4]{Avramov}.  Hence
\[
 H^q(\mathfrak g)=T^q_0(A).                                  \tag{2.5}
\]
The identical argument with $\Omega_{P/S}$ gives
\[
 \Der_S(P,P)_0\longrightarrow\Der_S(P,A)_0
\quad\text{as a quasi-isomorphism}.                          \tag{2.6}
\]
Here $\Der_S^0(P,A)=0$, and a relative degree-one cocycle is exactly an
$S$-linear map $I\to A$ of internal degree zero.  Consequently
$H^1(\mathfrak h)=\operatorname{Hom}_S(I,A)_0$.  The exact computation in
\path{verify_relative_controller.m2} obtains dimensions $109$ for this
space, $56$ for the graded Jacobian/coordinate-boundary image, and $53$ for
the absolute quotient.
The bracket used below is the bracket induced by the literal commutator in
(2.1), not a product on a conormal Ext group.

\subsection{Maurer--Cartan elements give embedded flat algebras}

For $R_n$, let
\[
 \widehat{\mathfrak h}_n
 =\mathfrak h\otimes_k(s)/(s^n),
\]
and over $R$ let
\[
 \widehat{\mathfrak h}=s\mathfrak h[[s]],
 \qquad F^r\widehat{\mathfrak h}=s^r\mathfrak h[[s]].       \tag{2.7}
\]
This filtration is complete and separated, the differential preserves it,
and brackets add filtration orders.

\begin{lemma}[filtered lifting of a Tate differential]\label{lem:cyclelift}
Let $P_{<q}$ be the subalgebra generated by the Tate generators of
homological degree less than $q$, and assume
$H_{q-2}(P_{<q},d)=0$, where $q\geq3$.  Let $D$ be a square-zero
weight-preserving differential on $P_{<q}\otimes R_n$ with
$D\equiv d\pmod s$.  Every homogeneous $d$-cycle
$c_0\in(P_{<q})_{q-1}$ has a homogeneous lift
$c\equiv c_0\pmod s$ satisfying $Dc=0$.
\end{lemma}

\begin{proof}
Construct
\(
 c^{(r)}=c_0+su_1+\cdots+s^{r-1}u_{r-1}
\)
so that $Dc^{(r)}\in s^rP_{<q}\otimes R_n$.  This holds for $r=1$
because $dc_0=0$.  Write
\[
 Dc^{(r)}=s^r\rho_r\pmod {s^{r+1}}.
\]
The identity $D^2=0$ on $P_{<q}$ gives
\[
 0=D(Dc^{(r)})=s^rd\rho_r\pmod {s^{r+1}},
\]
so $d\rho_r=0$.  The class of $\rho_r$ lies in homological degree $q-2$;
by hypothesis choose a homogeneous $u_r$ with $du_r=-\rho_r$.  Then
\[
 c^{(r+1)}=c^{(r)}+s^ru_r
\]
satisfies $Dc^{(r+1)}\in s^{r+1}P_{<q}\otimes R_n$.  After the
$(n-1)$st step its image $c$ in $R_n$ is a $D$-cycle.  Taking the internal
homogeneous component matching $c_0$ at every step is legitimate because
$d$ and $D$ preserve weight.
\end{proof}

\begin{proposition}[direct classical realization]\label{prop:realization}
Let $\alpha\in\widehat{\mathfrak h}_n^1$ satisfy
\[
 \partial\alpha+\frac12[\alpha,\alpha]=0.                   \tag{2.8}
\]
Then $D=d+\alpha$ is a square-zero, $S\otimes R_n$-linear,
weight-preserving derivation on $P\otimes R_n$.  If
\[
 B_n=H_0(P\otimes R_n,D),
\]
then $B_n$ is a flat graded embedded deformation of $A$ in the fixed
coordinates.  If $e_1,\ldots,e_{16}$ are the degree-one Tate generators with
$d(e_i)=m_i$, then, literally,
\[
 B_n=R_n[x_1,\ldots,x_8]/(D(e_1),\ldots,D(e_{16})),           \tag{2.9}
\]
and every $D(e_i)$ is a cubic.

Conversely, every marked flat graded quotient of $R_n[x_1,\ldots,x_8]$
specializing to $A$ admits, after choices of homogeneous Tate lifts, a
differential $D=d+\alpha$ of this form.  Conjugating $D$ by a nilpotent
relative degree-zero automorphism changes the semi-free trivialization while
fixing $S$ and induces an isomorphic embedded deformation.  No converse
classification of all isomorphisms by gauge is needed below.
\end{proposition}

\begin{proof}
For a degree-one derivation,
\[
 (d+\alpha)^2=\partial\alpha+\frac12[\alpha,\alpha],
\]
so (2.8) is exactly $D^2=0$.  Relative derivations vanish on $S$, and weight
zero makes $D(e_i)$ homogeneous of internal degree three.  Since every
homological-degree-one element is an $S\otimes R_n$-linear combination of
the $e_i$, the degree-zero boundaries form exactly the ideal in (2.9).

Filter $(P\otimes R_n,D)$ by powers of $s$.  The filtration is finite, and
its associated-graded differential is $d$.  Its spectral sequence has
\[
 E_1=H(P)\otimes_k\gr(R_n)
     =A\otimes_k\gr(R_n)
\]
in homological degree zero and is zero in positive homological degrees.
Thus $(P\otimes R_n,D)$ is exact in positive degrees.  In each internal
degree it is a free $R_n$-resolution of the finite module $(B_n)_j$.
Tensoring this resolution with $k$ gives $P_j$, so
\[
 \Tor^{R_n}_1(k,(B_n)_j)=H_1(P_j)=0.
\]
The local criterion for flatness~\cite[Tag 00MK]{StacksFlat} makes every
$(B_n)_j$ flat, hence free, over the Artinian local ring $R_n$.

For the converse, let $J$ be the kernel of a marked flat quotient
$R_n[x]\twoheadrightarrow B_n$.  Degreewise base change in
\[
 0\longrightarrow J\longrightarrow R_n[x]\longrightarrow B_n
 \longrightarrow0
\]
is exact.  Lift the basis $m_1,\ldots,m_{16}$ of $I_3$ to elements
$F_i\in J_3$.  Nakayama's lemma, first in degree three and then in every
degree, says that the $F_i$ generate $J$.  The kernel of
\[
 R_n[x](-3)^{16}\xrightarrow{(F_1,\ldots,F_{16})}R_n[x]
\]
also commutes with reduction, as follows in two steps.  In
$0\to J\to R_n[x]\to B_n\to0$, both the middle and right modules are
$R_n$-flat, so $J$ is $R_n$-flat.  If $K$ is the displayed first-syzygy
kernel, then
$0\to K\to R_n[x](-3)^{16}\to J\to0$ and
$\Tor^{R_n}_1(k,J)=0$ identify $K\otimes k$ with the full special kernel.
Hence the
thirty special first syzygies lift to genuine relations among the $F_i$.
Define $D$ on $e_i$ and $f_j$ by these equations and lifted relations.

It remains only to extend $D$ over the later Tate generators.  If $z$ has
homological degree $q\geq3$, then $dz$ is a cycle of degree $q-1$ in
$P_{<q}$.  By the staged Tate construction, all homology through degree
$q-2$ has already been killed.  Lemma~\ref{lem:cyclelift} therefore gives a
$D$-cycle $c_z\equiv dz\pmod s$; set $D(z)=c_z$.  Induction in the Tate
generators constructs $D$.  Since the square of an odd derivation is again a
derivation, checking $D^2$ on the generators proves $D^2=0$ everywhere.  The
difference $D-d$ is the required relative Maurer--Cartan element.  A
nilpotent exponential of a relative degree-zero derivation conjugates $D$ and
gives the stated one-way gauge assertion.
\end{proof}

The converse is used below in an even stronger form: the relevant
$R_4$ differential is constructed explicitly through the first relations,
and the remaining extension is the last paragraph of the proof.  Thus no
abstract comparison of deformation hulls is invoked.

\section{The exact \texorpdfstring{$R_4$}{R4} deformation is an actual
Maurer--Cartan element}

\subsection{The certified presentation and flatness}

The exact tangent direction in the ordered $53$-element absolute tangent
basis is
\[
\begin{split}
y=(\,&1,3,1,2,6,1,1,1,1,1,21,28,1,5,10,55,66,15,1,18,2,\\
&35,42,1,4,10,1,25,12,30,5,33,44,1,6,3,14,7,1,9,1,6,\\
&12,20,1,4,22,1,15,1,1,8,11\,).
\end{split}                                                    \tag{3.1}
\]

The pinned calculation produces matrices $F_0,F_1,F_2,F_3$ and
$Q_0,Q_1,Q_2,Q_3$.  After substituting $t=sy$, set
\[
 F^{[3]}=\sum_{i=0}^3F_i,
 \qquad Q^{[3]}=\sum_{i=0}^3Q_i
\]
over $R_4[x_1,\ldots,x_8]$.

\begin{computation}[starting deformation]\label{comp:start}
The script \path{../completion/verify_starting_jet.m2} verifies, by exact
arithmetic over $\QQ$, that
\[
 F^{[3]}Q^{[3]}=0,                                           \tag{3.2}
\]
that $Q_0$ is the complete thirty-column first-syzygy matrix of $F_0$, and
that $F^{[3]}\bmod s^3$ is byte-for-byte the sixteen-row file
\path{../referee-packet/data/universal_2jet_QQ.txt}.  Its SHA--256 is
\[
\texttt{40e64e61674b6a4e61f1ea6822dc79327bf4ba397285f84ed8738c4c18cd1795}.
\tag{3.3}
\]
The certificate is \path{../completion/verification/starting_jet_QQ.txt}.
\end{computation}

Let
\[
 A_4=R_4[x]/(F^{[3]}).
\]
This is an actual flat algebra, not merely a matrix satisfying a formal
equation.  In an internal degree $j$, let $K_j$ be the kernel of the map
defined by $F^{[3]}$.  Since the columns of $Q^{[3]}$ lie in $K_j$ and reduce
to a generating set of the entire special kernel, the reduction map
\[
 K_j\otimes_{R_4}k\longrightarrow\ker(F_0)_j
\]
is onto.  Extend the presentation on the left by a free module.  After
tensoring with $k$, this surjectivity says precisely that
\(
 \Tor^{R_4}_1(k,(A_4)_j)=0.
\)
The local criterion~\cite[Tag 00MK]{StacksFlat} makes $(A_4)_j$ free.  This
holds in every internal degree, so $A_4$ is a flat graded embedded
deformation.  Its quotient $A_4/s^3A_4$ is the printed two-jet.

\subsection{A differential on the audited Tate resolution}

We now construct the promised Maurer--Cartan representative rather than
deducing one from the package terminology.  The certified beginning of $P$
has generators
\[
\begin{array}{c|c|c}
 &\text{homological degree}&\text{internal weight}\\ \hline
e_1,\ldots,e_{16}&1&3\\
f_1,\ldots,f_{30}&2&4\\
g_1,\ldots,g_{136}&3&5\text{ (16 times), }6\text{ (120 times).}
\end{array}                                                     \tag{3.4}
\]
Let $R_0$ be the relation matrix used for $d(f)$ in this Tate model.  The
strict bracket script computes an invertible constant matrix $U$ with
\[
 Q_0U=R_0,
 \qquad \det U=1.                                             \tag{3.5}
\]
On $P\otimes R_4$, define
\[
 D(e_i)=F_i^{[3]},
 \qquad
 D(f_j)=\sum_i(Q^{[3]}U)_{ij}e_i.                             \tag{3.6}
\]
Equations (3.2) and (3.5) give $D^2(e_i)=D^2(f_j)=0$, and
$D\equiv d\pmod s$.

Extend $D$ over each $g_a$ and then every later Tate generator by
Lemma~\ref{lem:cyclelift}.  For a generator of homological degree $q\geq3$,
the obstruction at each $s$-order is a $d$-cycle in degree $q-2>0$ and has a
homogeneous primitive in degree $q-1$, which cannot involve that generator or
a later one.  The
completeness of the thirty $f$-stage and the $136$-generator $g$-stage needed
here is independently certified in Section~\ref{sec:h3}; all later stages are
Tate stages by construction.  We obtain a square-zero differential on the
same underlying semi-free algebra,
\[
 D=d+\bar\alpha,
 \qquad
 \bar\alpha\in
 \MC\bigl(\mathfrak h\otimes(s)/(s^4)\bigr).                 \tag{3.7}
\]
Its degree-zero homology is exactly $A_4$, because its degree-one boundaries
are the sixteen entries of $F^{[3]}$.  Its linear coefficient is the strict
degree-one cocycle whose values on the $e_i$ are $T^1y$; its image in
$H^1(\mathfrak g)$ is therefore the computed class $y$.  This proves the
required connection between the $R_4$ matrices and an actual truncated
Maurer--Cartan element.

\section{The certified degree-\texorpdfstring{$1,2,3$}{1,2,3} exact
sequence}\label{sec:h3}

\subsection{What the 136 cycles certify}

Let $D_1,D_2$ be the chain maps in homological degrees one and two of the
materialized Tate algebra.  Let $D_3$ denote the differentials of the
decomposable degree-three elements $e_if_j$ and $e_ie_je_k$, and let $Z$ be
the $150\times136$ matrix whose columns are the proposed $d(g_a)$.

\begin{computation}[complete low Tate stages]\label{comp:tate}
The independent audit in \path{../audits/tate-stage/} reconstructs all four
matrices from the canonical sparse input and proves over $\QQ[x_1,\ldots,x_8]$
\[
 D_1D_2=0,\qquad D_2D_3=0,\qquad D_2Z=0,
\]
and the unbounded graded-module equalities
\[
 \ker D_1=\im D_2,
 \qquad
 \ker D_2=\im[D_3\;Z].                                      \tag{4.1}
\]
Macaulay2 computes full homogeneous syzygy modules and two-sided lifts;
Singular independently recomputes the modules and lifts.  The result is
\path{../audits/tate-stage/MATHEMATICAL_CERTIFICATE.md}.
\end{computation}

The precise meaning of (4.1) is worth emphasizing.  The $136$ columns are
cycles in homological degree two.  Their classes span
\[
 \ker D_2/\im D_3,                                           \tag{4.2}
\]
so adjoining the $g_a$ kills all remaining degree-two homology.  They are not
asserted to be a basis of $H^3$, and no total dimension of $H^3$ is used.
Equation (4.1) proves that no additional Tate generators of homological
degree at most three are missing; a full Tate resolution can continue with
generators of degree at least four.

\subsection{Why the computed columns are genuine
\texorpdfstring{$H^3$}{H3} classes}

Let
\[
 C^p=\Der_k^p(P,A)_0.
\]
On (3.4),
\[
 C^1=(A_3)^{16},\quad C^2=(A_4)^{30},\quad
 C^3=(A_5)^{16}\oplus(A_6)^{120}.                            \tag{4.3}
\]
The exact bracket calculation produces a closed degree-one derivation
$\theta$ representing $y$ and closed degree-two derivations
$\eta_1,\ldots,\eta_{27}$ representing a basis of $H^2$.  It solves and
rechecks their closure equations on every $e_i,f_j,g_a$.
It also constructs degree-one derivations
$\psi_1,\ldots,\psi_{53}$ for the absolute tangent basis and checks their
closure through every $f_j$; their later extension is included in the next
lemma.

\begin{lemma}[extension of the computed derivations]\label{lem:derextend}
Each of $\theta,\eta_1,\ldots,\eta_{27}$ extends, without changing its values
on (3.4), to a closed derivation of the full Tate resolution.  Each
$\psi_i$ extends without changing its values on the $e$- and $f$-generators.
\end{lemma}

\begin{proof}
For a $\psi_i$, the first unchecked generators are the $g_a$ of degree three.
The right side of its closure equation is a degree-one cycle.  The equality
$\ker D_1=\im D_2$ in Computation~\ref{comp:tate} says exactly that this
cycle has a degree-two primitive in the algebra generated by the $e$ and
$f$.  Thus every $\psi_i$ extends through all the $g_a$ without changing the
values used by the primary-commutator calculation.

Suppose a degree-$p$ derivation $\vartheta$, with $p=1$ or $2$, is already
closed on earlier generators, and let $z$ be a new Tate generator of degree
$q\geq4$.  Closedness on $z$ requires
\[
 d(\vartheta z)=(-1)^p\vartheta(dz).                          \tag{4.4}
\]
The right side is a cycle of homological degree $q-p-1>0$.  It is a boundary
because $P$ is acyclic in positive degree.  A primitive has degree
$q-p<q$, hence uses only earlier generators.  Its component of the required
internal weight is still a primitive.  Induction proves the assertion.  The
low cases excluded by the inequality are exactly the closure equations
checked on (3.4).
\end{proof}

The commutators $[\theta,\eta_j]$ are therefore global closed degree-three
derivations.  Under (2.4), their cochains lie in the kernel of the outgoing
map $C^3\to C^4$.  Consequently
\[
 H^3(C)=\ker(C^3\to C^4)/\im(C^2\to C^3)
 \hookrightarrow C^3/\im(C^2\to C^3).                       \tag{4.5}
\]
Independence of these already-closed columns in the quotient on the right is
therefore independence in genuine $H^3$.  This is the rigorous role of the
$136$-cycle certificate; a bare degree-three generator count would not prove
(4.5).

\subsection{Exactness, with both maps computed}

\begin{proposition}[fixed-direction strict exactness]\label{prop:exact}
Let $a\in Z^1(\mathfrak h)$ be the linear coefficient of (3.7).  Then
\[
 H^1(\mathfrak h)\xrightarrow{\ad_a}H^2(\mathfrak h)
 \xrightarrow{\ad_a}H^3(\mathfrak h)                        \tag{4.6}
\]
is exact at $H^2$.  The image of the first map has dimension $15$ and the
second map has rank $12$ on the $27$-dimensional middle space.
\end{proposition}

\begin{proof}
The script \path{../referee-packet/code/verify_aq_bracket.m2} performs the
following exact operations.

First relate its chosen cocycle $\theta$ to the coefficient $a$ of the
actual differential (3.7).  Both have absolute cohomology class $y$, by the
last paragraph of Section~3, so
\[
 a-\theta=\partial\xi\qquad\text{for some }\xi\in\mathfrak g^0.
 \tag{4.7}
\]
If $z$ is any closed relative derivation of positive degree, then
\[
 [a-\theta,z]=[\partial\xi,z]=\partial[\xi,z].              \tag{4.8}
\]
The cochain $[\xi,z]$ has positive degree and is therefore relative by
(2.2).  Thus $\ad_a$ and $\ad_\theta$ induce the same maps on relative
$H^1$ and $H^2$.  In particular, it is enough to compute with the literal
$\theta$ used by the script; this is a chain-level comparison, not an appeal
to invariance of any transferred operation.

First, it transports all $27$ package obstruction columns through the
constant relation-basis change (3.5), checks that they are closed, proves
that they are independent modulo boundaries, and proves that no further
internal-degree-zero $H^2$ cocycle remains.  Thus they are a basis of the
$27$-dimensional $H^2$.

Second, it computes the literal graded commutators
\[
 [\theta,\eta_j]=\theta\eta_j-\eta_j\theta
\]
on all $136$ generators $g_a$.  By Lemma~\ref{lem:derextend} and (4.5), the
rank $12$ found modulo every degree-three boundary is the rank of
\[
 A_y=\ad_y:H^2(\mathfrak g)\longrightarrow H^3(\mathfrak g).
\]

Third, it separately computes the primary commutators of $\theta$ with all
$53$ degree-one absolute tangent basis derivations.  For two degree-one
derivations the bracket is the signed sum
$\theta\psi+\psi\theta$.  The resulting map
\[
 B_y=\ad_y:H^1(\mathfrak g)\longrightarrow H^2(\mathfrak g)
\]
has rank $15$.  More precisely, the script proves the matrix identity
\[
 B_y=-J\,Dq_y,                                                \tag{4.9}
\]
where $J$ is the verified isomorphism from the $27$ package obstruction
coordinates to the displayed $H^2$ basis.  It also checks directly, on those
strict commutator matrices,
\[
 A_yB_y=0.                                                    \tag{4.10}
\]
Thus (4.10), not the rank of $Dq_y$ alone, gives
$\im B_y\subseteq\ker A_y$.  Rank--nullity gives
\[
 \dim\ker A_y=27-12=15=\dim\im B_y,
\]
so equality follows.

Finally transfer this statement from the absolute algebra to the relative
controller.  By (2.2)--(2.3), $H^2$ and $H^3$ and the second map are literally
the same, while every absolute degree-one class has a relative cocycle
representative.  Conversely, the extra relative classes become absolute
boundaries.  If such a class is represented by $\partial\xi$, with
$\xi\in\mathfrak g^0$, then
\[
 [a,\partial\xi]=-\partial[a,\xi],                            \tag{4.11}
\]
because $\partial a=0$; moreover $[a,\xi]$ has positive degree and is
therefore relative.  Hence the extra classes contribute zero to $H^2$, and
the images of the relative and absolute first maps coincide.  Exactness of
(4.6) follows.
\end{proof}

The deterministic output
\path{../referee-packet/verification/aq_bracket_QQ.txt} records the $27$
classes, ranks $12$ and $15$, the zero composite, and the primary identity
(4.9).  Notice that the proof neither asserts nor needs a dimension for the
whole group $H^3$.

\section{Strict Bianchi induction to all orders}

For $\gamma\in s\mathfrak h^1[[s]]$, define its curvature by
\[
 \mathcal F(\gamma)=\partial\gamma+\frac12[\gamma,\gamma].   \tag{5.1}
\]
The dg Lie axioms give the strict Bianchi identity
\[
 \partial\mathcal F(\gamma)+[\gamma,\mathcal F(\gamma)]=0.  \tag{5.2}
\]
Indeed,
$\partial(\frac12[\gamma,\gamma])=-[\gamma,\partial\gamma]$
and graded Jacobi gives $[\gamma,[\gamma,\gamma]]=0$.

\begin{lemma}[fixed-two-jet extension]\label{lem:induction}
Let $\bar\alpha\in\MC(\mathfrak h\otimes(s)/(s^4))$ have linear
coefficient $a$.  If (4.6) is exact at $H^2$, then there is
\[
 \alpha_\infty\in\MC(s\mathfrak h[[s]])
 \quad\text{with}\quad
 \alpha_\infty\equiv\bar\alpha\pmod{s^3}.                  \tag{5.3}
\]
\end{lemma}

\begin{proof}
Lift $\bar\alpha$ arbitrarily to $\gamma=sa+O(s^2)$.  Its curvature is
$O(s^4)$.  Suppose, for $N\geq4$, that
\[
 \mathcal F(\gamma)=s^Nr_N+s^{N+1}r_{N+1}+O(s^{N+2}).       \tag{5.4}
\]
The coefficients of $s^N$ and $s^{N+1}$ in (5.2) give
\[
 \partial r_N=0,
 \qquad \partial r_{N+1}+[a,r_N]=0.                         \tag{5.5}
\]
Thus the obstruction coefficient is a cocycle and its class lies in the
kernel of $\ad_a:H^2\to H^3$.  Exactness supplies a closed
$v\in\mathfrak h^1$ such that
\[
 [r_N]+[a,v]=0.
\]
Choose $w\in\mathfrak h^1$ with
\[
 r_N+[a,v]=\partial w.                                      \tag{5.6}
\]
Replace
\[
 \gamma\quad\text{by}\quad
 \gamma'=\gamma+s^{N-1}v-s^Nw.                             \tag{5.7}
\]
The exact change formula
\[
 \mathcal F(\gamma')-\mathcal F(\gamma)
 =\partial(\gamma'-\gamma)+[\gamma,\gamma'-\gamma]
  +\frac12[\gamma'-\gamma,\gamma'-\gamma]
\]
shows that the order-$s^N$ change is $[a,v]-\partial w$, which cancels $r_N$
by (5.6).  All other terms have order at least $N+1$: $\partial v=0$,
$\gamma-sa=O(s^2)$, and $2N-2\geq N+1$.

The corrections are Cauchy in the filtration (2.7), so they converge
coefficientwise.  Continuity of (5.1), together with
$\bigcap_Ns^N\mathfrak h[[s]]=0$, gives zero curvature in the limit.  At the
initial $N=4$ step the correction is $s^3v-s^4w$; consequently the
coefficients of $s$ and $s^2$ are unchanged, proving (5.3).
\end{proof}

Apply the lemma to (3.7) and Proposition~\ref{prop:exact}.  This gives
$\alpha_\infty$.  It is important that the conclusion preserves the
$R_3$ reduction, not necessarily the whole $R_4$ element: the first
correction may change the coefficient of $s^3$.  Preserving the printed
first and second-order terms is exactly what the smoothness certificate
requires.

\section{The formal embedded family and its algebraization}

Put $D_\infty=d+\alpha_\infty$ and define the literal cubics
\[
 F_i(s)=D_\infty(e_i)\in R[x_1,\ldots,x_8]_3.                \tag{6.1}
\]
By (5.3) and Computation~\ref{comp:start},
\[
 F_i(s)=F_i^{(2)}+s^3H_i(s,x),                               \tag{6.2}
\]
where the sixteen $F_i^{(2)}$ are exactly the rows of the frozen two-jet and
$H_i$ is homogeneous cubic in $x$.  There are only finitely many cubic
monomials, so (6.1) is an honest element of $R[x]$, not a second completion
in the $x$-variables.

Let
\[
 B=R[x_1,\ldots,x_8]/(F_1(s),\ldots,F_{16}(s)),
 \qquad X=\Proj B.                                           \tag{6.3}
\]
For every $n$, Proposition~\ref{prop:realization} applied to
$\alpha_\infty\bmod s^n$ gives
\[
 B/s^nB
 =H_0(P\otimes R_n,D_\infty\bmod s^n),                     \tag{6.4}
\]
and every internal-degree piece of (6.4) is free over $R_n$.

We check flatness over $R$, rather than assuming it from formal effectivity.
Each $B_j$ is a finite $R$-module.  If $sb=0$ in $B_j$, then flatness of
$B_j/s^nB_j$ over $R_n$ implies
\[
 b\in s^{n-1}B_j\quad\text{for every }n\geq2,
\]
because the annihilator of $s$ on a free $R_n$-module is $s^{n-1}$ times
that module.  Krull intersection gives $b=0$.  Hence $B_j$ is torsion-free,
and therefore free over the DVR $R$.  The direct sum $B=\bigoplus_jB_j$ is
$R$-flat.  On a standard projective chart, $(B_{x_i})_0$ is a direct summand
of the flat localization $B_{x_i}$, so $X\to\Spec R$ is flat.

Equation (6.3) is also an explicit algebraization of the compatible formal
family (6.4): it is a closed subscheme of $\PP^7_R$, is projective over $R$,
and has exactly those reductions.  Alternatively, the compatible closed
subschemes satisfy every hypothesis of the Stacks Project algebraization
lemma~\cite[Tag 0899]{StacksAlg}: $R$ is Noetherian and $(s)$-adically
complete, $\PP^7_R$ is separated of finite type, the squares are cartesian,
and $X_0$ is proper.  The explicit cubics identify that algebraized closed
subscheme with (6.3).  Thus no formal family is being silently substituted
for an algebraic one.

\section{The 280 incidence certificates and smoothness}

On the projective chart $x_v=1$, write $z_1,\ldots,z_7$ for the other affine
coordinates and let $J$ be the $7\times16$ affine Jacobian of the cubics.  A
four-plane in $\ker J^{\mathsf T}$ is represented on one of the $35$
standard charts of $\operatorname{Gr}(4,7)$ by a $7\times4$ matrix $K_P$
whose pivot rows form the identity.  For each of the $8\cdot35=280$ pairs,
the truncated singular-incidence ideal has the $80$ displayed generators
\[
 \mathcal J_{v,P}^{(2)}
 =(F_1^{(2)},\ldots,F_{16}^{(2)},(J^{(2)})^{\mathsf T}K_P).
\tag{7.1}
\]

\begin{computation}[all incidence charts]\label{comp:incidence}
The scripts \path{../referee-packet/code/verify_incidence_chart.m2} and
\path{../referee-packet/code/verify_all_incidence.py} prove, over exact
$\QQ$ arithmetic, that on every chart
\[
 s^2\in(\mathcal J_{v,P}^{(2)},s^3).                          \tag{7.2}
\]
They encode the $80$ generators of (7.1) over
$\QQ[z_1,\ldots,z_7,a_1,\ldots,a_{12}][s]/(s^3)$ as a
$3\times240$ module matrix and prove that $(0,0,1)^{\mathsf T}$ belongs to
its column module.  Exact partial Gr\"obner reduction proves $277$ charts at
degree limit six; degree limit eight proves precisely $(v,P)=(2,7),(6,12)$,
and $(8,16)$.  The driver verifies that all $280$ pairs occur and checks one
retained lift back to the original columns.  The certificate is
\path{../referee-packet/verification/incidence_all_QQ.txt}.
\end{computation}

The certificates assert exactly (7.2).  They do not, by themselves, assert
that a fibre is smooth; that conclusion is the following geometric argument.

\begin{proposition}[two-jet smoothness]\label{prop:smooth}
Every flat projective continuation satisfying (6.2) has smooth geometric
generic fibre.
\end{proposition}

\begin{proof}
The special fibre is Cohen--Macaulay and equidimensional of dimension three
by Computation~\ref{comp:special}.  For a proper flat family over $R$, the
Cohen--Macaulay fibre criterion and constancy of relative dimension make the
morphism Cohen--Macaulay of relative dimension three; these are exactly the
opening reductions in the Stacks smoothing lemma~\cite[Tag 0E7T]{StacksSmooth}.
Thus a point of a fibre is singular precisely when the affine Jacobian has
rank at most three, equivalently when $\ker J^{\mathsf T}$ contains a
four-plane.

Suppose the geometric generic fibre has a singular point.  The incidence
scheme is of finite type over $K=\QQ((s))$, so it has a closed point after a
finite extension $K'/K$.  Let $R'$ be a DVR of the integral closure of $R$
in $K'$ lying over $(s)$.  Properness of $\PP^7_R$ extends the point over
$R'$, and properness of $\operatorname{Gr}(4,7)_R$ extends a chosen
four-plane in the kernel.  After scaling, some projective coordinate and some
Pl\"ucker coordinate are units.  Hence one of the $280$ chart pairs has all
of its point and four-plane coordinates in $R'$.

On this chart, (7.2) is an exact polynomial identity
\[
 s^2=\sum_jc_jq_j+s^3c_0,                                   \tag{7.3}
\]
where the $q_j$ are the truncated incidence generators.  Equation (6.2)
changes each equation, and each affine partial derivative, only by $s^3$
times an integral power series.  Replacing $q_j$ by the actual incidence
generator $q'_j$ in (7.3) therefore gives
\[
 s^2=\sum_jc_jq'_j+s^3c
\]
with $c\in R'$ after evaluation at the integral incidence section.  All
$q'_j$ vanish there, so $s^2=s^3c$.  The valuation of $s$ in $R'$ is
positive, whereas this equality would give
\[
 2v(s)=3v(s)+v(c)\geq3v(s),
\]
a contradiction.  The geometric generic fibre is smooth.
\end{proof}

\begin{proof}[Proof of Theorem~\ref{thm:main}]
Computation~\ref{comp:start} and the direct Tate lift (3.6)--(3.7) give a
truncated relative Maurer--Cartan element with the certified two-jet.
Proposition~\ref{prop:exact} and Lemma~\ref{lem:induction} produce the
all-orders strict Maurer--Cartan element without changing that two-jet.
Equations (6.1)--(6.4) give a flat projective family with special fibre
$X_0$, and Proposition~\ref{prop:smooth} proves that its geometric generic
fibre is smooth.
\end{proof}

\section{Logical and computational dependency ledger}

\begin{center}
\begin{tabular}{p{0.19\textwidth}p{0.34\textwidth}p{0.35\textwidth}}
\toprule
Statement & Evidence & Exact consequence used\\
\midrule
Special fibre & Exact combinatorial replay; Reisner's criterion & Pure CM
threefold with ideal (1.1)\\
$R_4$ start & Exact $F^{[3]}Q^{[3]}=0$, complete special syzygies, byte
comparison & Flat embedded deformation and literal two-jet\\
Low Tate stage & Independent Macaulay2 and Singular full syzygy modules &
$136$ cycles span (4.2); no missing low generators\\
Strict exactness & Closed global derivations; all primary and secondary
commutators; ranks and direct zero composite & Exactness (4.6), without a
total $H^3$ dimension\\
All orders & Equations (5.1)--(5.7) & Complete relative Maurer--Cartan element
fixing coefficients of $s,s^2$\\
Generic smoothness & $280$ exact memberships plus projective/Grassmannian
properness & Valuation contradiction for every higher continuation\\
Algebraization & Sixteen coefficientwise cubic series; alternatively Tag
0899 & Actual projective $R$-scheme, not only formal data\\
\bottomrule
\end{tabular}
\end{center}

The computations do not prove a total value for $\dim H^3$, do not preserve
the full $s^3$ coefficient of the starting $R_4$ object, and do not identify
the generic fibre beyond its smoothness.  None of those stronger statements
is needed.

\begin{thebibliography}{9}

\bibitem{Avramov}
L.~L.~Avramov,
\emph{Infinite free resolutions}, in \emph{Six Lectures on Commutative
Algebra}, Progress in Mathematics 166, Birkh\"auser, 1998, pp.~1--118.

\bibitem{GS}
B.~Gr\"unbaum and V.~P.~Sreedharan,
\emph{An enumeration of simplicial $4$-polytopes with $8$ vertices},
J. Combinatorial Theory \textbf{2} (1967), 437--465,
\href{https://doi.org/10.1016/S0021-9800(67)80055-3}
{doi:10.1016/S0021-9800(67)80055-3}.

\bibitem{Reisner}
G.~A.~Reisner,
\emph{Cohen--Macaulay quotients of polynomial rings},
Adv. Math. \textbf{21} (1976), 30--49,
\href{https://doi.org/10.1016/0001-8708(76)90114-6}
{doi:10.1016/0001-8708(76)90114-6}.

\bibitem{StacksFlat}
The Stacks Project Authors,
\emph{The Stacks Project},
\href{https://stacks.math.columbia.edu/tag/00MK}{Tag 00MK}, local criterion
for flatness.

\bibitem{StacksAlg}
The Stacks Project Authors,
\emph{The Stacks Project},
\href{https://stacks.math.columbia.edu/tag/0899}{Tag 0899}, Grothendieck
algebraization for compatible closed subschemes.

\bibitem{StacksSmooth}
The Stacks Project Authors,
\emph{The Stacks Project},
\href{https://stacks.math.columbia.edu/tag/0E7T}{Tag 0E7T}, smoothing lemma.

\end{thebibliography}

\end{document}
